% context: 9th-grade physics classwork, applications of Unit 1 skills
% topic: unit conversion chains, scientific notation, significant figures,
%        estimation (adapted from Giancoli Ch. 1 problems 49-51)
% structure: two pages, four problems, widely spaced for written work
\documentclass[12pt, twoside]{article}
\input{../preamble}
\usepackage{float}

\title{Physics}
\author{Chris Huson}
\date{September 2026}

\fancyhead[LE]{\thepage}
\fancyhead[RO]{\thepage \\ Nickname: \hspace{2.25cm} \,\\ \,}
\fancyhead[LO]{La Scuola d'Italia / Huson / Physics \\* 25 September 2026}

\begin{document}

\subsubsection*{1.5 Classwork: Applications}

\textbf{Reference (provided).} \quad
average speed $= \dfrac{\text{distance}}{\text{time}}$ \quad
area of a circle $= \pi r^2$ \\[0.1cm]
speed of light $c = 3.00 \times 10^{8}$ m/s \quad 1 km = 1000 m \quad 1 m = 100 cm \\
1 h = 3600 s \quad 1 min = 60 s \quad 1 yr $\approx$ 365 days \quad
1 metric ton = 1000 kg \\
1 m$^3$ = 1000 L \quad 1 L of water has a mass of 1 kg \quad 1 gal $\approx$ 3.785 L

\vspace{0.4cm}

\begin{enumerate}[itemsep=0.5cm]

\item During a heavy storm, 2.5 cm of rain falls on Manhattan. Treat Manhattan as a
  rectangle 20 km long and 3.0 km wide.
  \begin{enumerate}[itemsep=4.6cm]
    \item What volume of water fell, in cubic meters? Convert each length to meters
      first.
    \item What is the mass of this water, in kilograms and in metric tons?
    \item How many gallons is it? New York City uses about 1 billion gallons of water
      a day --- how does the storm compare?
  \end{enumerate}

\newpage

\item In 1773 Benjamin Franklin poured a teaspoon of oil on a pond and watched it
  spread into a thin film. Suppose 1.0 L of oil spreads over a calm lake until the
  slick is just one molecule thick, about $2 \times 10^{-10}$ m.
  \begin{enumerate}[itemsep=2.1cm]
    \item What is the volume of the oil in cubic meters?
    \item The volume of the slick equals its area times its thickness. What is the
      area of the slick, in square meters?
    \item If the slick is a circle, estimate its diameter.
  \end{enumerate} \vspace{2.0cm}

\item The distance around the Earth at the equator is 40{,}075 km. Suppose you could
  walk it at 5.0 km/h, walking 12 hours each day.
  \begin{enumerate}[itemsep=2.0cm]
    \item How many days would the trip take?
    \item About how many years is that?
  \end{enumerate} \vspace{2.0cm}

\item The Sun is $1.496 \times 10^{8}$ km from the Earth. How many minutes does it
  take sunlight to reach us? Write the conversion as a single chain of fractions.

\end{enumerate}
\end{document}
